\documentclass[11pt]{article}

\DeclareUnicodeCharacter{0301}{\'e}

\usepackage[margin=1in]{geometry}
\usepackage{multirow}
\usepackage{amssymb,amsmath,amsthm,amsfonts}
\usepackage{bm}
\usepackage{textgreek}
\usepackage{mathtools}
\usepackage{enumitem}
\usepackage[numbers,comma,sort&compress]{natbib}
\usepackage{authblk}
\usepackage{graphicx}
\usepackage[font=small]{caption}
\usepackage{subcaption} % [labelformat=simple]
\usepackage{tablefootnote} 
\usepackage{float}
\usepackage[ruled,vlined,linesnumbered]{algorithm2e}
\usepackage{algorithmic}
\renewcommand{\algorithmicrequire}{\textbf{Input: }}  
\renewcommand{\algorithmicensure}{\textbf{Output: }}
\usepackage{physics}
\usepackage{footnote}
\usepackage{xcolor}
\usepackage{mathrsfs}
\usepackage{bbm}
\usepackage{makecell}
\usepackage{pbox}
\usepackage[colorlinks]{hyperref}
\hypersetup{citecolor=blue}
% \urlstyle{same}
%\captionsetup[figure]{labelfont=bf}

\newtheorem{fact}{Fact}[section] 
\newtheorem{theorem}{Theorem}
\newtheorem{definition}{Definition}
\newtheorem{lemma}{Lemma}
\newtheorem{proposition}{Proposition}
\newtheorem{example}{Example}
\newtheorem{corollary}{Corollary}
\newtheorem{remark}{Remark}
\newtheorem{assumption}{Assumption}
\newcommand{\eqn}[1]{(\ref{eqn:#1})}
\newcommand{\eq}[1]{(\ref{eq:#1})}
\newcommand{\prb}[1]{(\ref{prb:#1})}
\newcommand{\thm}[1]{\hyperref[thm:#1]{Theorem~\ref*{thm:#1}}}
\newcommand{\cor}[1]{\hyperref[cor:#1]{Corollary~\ref*{cor:#1}}}
\newcommand{\defn}[1]{\hyperref[defn:#1]{Definition~\ref*{defn:#1}}}
\newcommand{\lem}[1]{\hyperref[lem:#1]{Lemma~\ref*{lem:#1}}}
\newcommand{\prop}[1]{\hyperref[prop:#1]{Proposition~\ref*{prop:#1}}}
\newcommand{\assum}[1]{\hyperref[assum:#1]{Assumption~\ref*{assum:#1}}}
\newcommand{\fig}[1]{\hyperref[fig:#1]{Figure~\ref*{fig:#1}}}
\newcommand{\tab}[1]{\hyperref[tab:#1]{Table~\ref*{tab:#1}}}
\newcommand{\algo}[1]{\hyperref[algo:#1]{Algorithm~\ref*{algo:#1}}}
\renewcommand{\sec}[1]{\hyperref[sec:#1]{Section~\ref*{sec:#1}}}
\newcommand{\append}[1]{\hyperref[append:#1]{Appendix~\ref*{append:#1}}}
\newcommand{\fac}[1]{\hyperref[fac:#1]{Fact~\ref*{fac:#1}}}
\newcommand{\lin}[1]{\hyperref[lin:#1]{Line~\ref*{lin:#1}}}
\newcommand{\itm}[1]{\ref{itm:#1}}
\newcommand{\bdelta}{\mathbf{\Delta}}
\renewcommand{\arraystretch}{1.25}
\newcommand{\algname}[1]{\textup{\texttt{#1}}}
\renewcommand{\i}{\mathrm{i}}
\renewcommand{\d}{\mathrm{d}}
\renewcommand{\L}{\mathcal{L}}
\newcommand{\ad}{\mathrm{ad}}
\newcommand{\T}{\mathcal{T}}
\newcommand{\B}{\Big}
\renewcommand{\P}{\mathcal{P}}
\newcommand{\simu}{\mathrm{sim}}
\newcommand{\Hc}{H_{\mathrm{control}}}
\def\>{\rangle}
\def\<{\langle}
\def\trans{^{\top}}

\newcommand{\xz}[1]{{\color{cyan} XZ: #1}}
\newcommand{\tyl}[1]{{\color{blue} Tongyang: #1}}
\newcommand{\zsy}[1]{{\color{teal} Shengyu: #1}}
\newcommand{\sz}[1]{{\color{purple} Shuo: #1}}



\title{Lindbladian Simulation with ZNE}
\author{}
\date{}

\begin{document}

\maketitle
\section{Error Bounds of the Richardson Extrapolation}
The dynamics of a digital quantum computer implementing a circuit under Markovian noise can be modeled using a master equation. The ideal evolution is driven by a time-dependent Hamiltonian $H(t)$ for $t\in[0,T]$ such that $\T \exp(-\i\int_{0}^T H(t)\,\d t)=U$. Concurrently, the system is subject to dissipative noise described by a Lindbladian $\lambda \L$. The dynamics of the system's density matrix $\rho(t)$ is described by \begin{align*}
    \frac{\d \rho(t)}{\d t} = -\i\ad_{H(t)}(\rho(t))+\lambda \L(\rho(t)),
\end{align*}
where $\ad_H(\rho) :=[H,\rho]$.

We define the ideal time-evolution operator and the noisy time-evolution operator as \begin{align*}
    \Phi(t_2,t_1):=\T \exp(\int_{t_1}^{t_2} -\i\ad_{H(t)}\,\d t),\quad \tilde{\Phi}_{\lambda}(t_2, t_1):=\T \exp(\int_{t_1}^{t_2}(-\i \ad_{H(t)}+\lambda \L)\,\d t)
\end{align*} Denote the initial state by $\rho_0$. By the variation of parameters formula, the noisy density matrix at time $T$ is \begin{align*}
    \tilde{\rho}(T) &= \tilde{\Phi}_{\lambda}(T,0)\rho_0\\
    &=\Phi(T,0)\rho_0+\B(\int_{0}^T \Phi(T, \tau_1)\lambda \L\tilde{\Phi}_{\lambda}(\tau_1, 0)\, \d \tau_1\B)\rho_0 \\
    & = U\rho_0U^{\dagger} + \B(\int_{0}^T \Phi(T, \tau_1)\lambda \L\Phi(\tau_1, 0)\, \d \tau_1\B)\rho_0 + \B(\int_{0}^T\d \tau_1 \int_{0}^{\tau_1}\d \tau_2 \Phi(t, \tau_1)\lambda \L\Phi(\tau_1, \tau_2)\lambda \L\tilde{\Phi}_{\lambda}(\tau_2,0)\B)\rho_0\\
    &\cdots \\
    &=U\rho_0U^{\dagger}+
    \sum_{\ell=1}^{K-1}\B(\int_{0}^T \d \tau_1 \int_{0}^{\tau_1}\d \tau_2 \cdots \int_{0}^{\tau_{\ell-1}} \d \tau_{\ell} \Phi(T, \tau_1)\lambda \L \Phi(\tau_1, \tau_2)\lambda \L \cdots \Phi(\tau_{\ell-1}, \tau_{\ell})\lambda \L\Phi(\tau_{\ell},0)\B)\rho_0\\
    &\quad+\B(\int_{0}^T \d \tau_1 \int_{0}^{\tau_1}\d \tau_2 \cdots \int_{0}^{\tau_{K-1}} \d \tau_{K} \Phi(T, \tau_1)\lambda \L \Phi(\tau_1, \tau_2)\lambda \L \cdots \Phi(\tau_{K-1}, \tau_{K})\lambda \L\tilde{\Phi}_{\lambda}(\tau_{K},0)\B)\rho_0.
\end{align*}
For any observable $O$, the expected value on $\tilde{\rho}(T)$ is \begin{align}
\label{eq:expansion-lambda}
    \tr[O\tilde{\rho}(T)] &= \tr[OU\rho_0 U^{\dagger}]+
    \sum_{\ell=1}^{K-1}F_{T,\ell}(O)\lambda^{\ell}+F_{T,K}(O,\lambda)\lambda^K,
\end{align}
where \begin{align*}
    F_{T,\ell}(O):=\tr\B[O\B(\int_{0}^T \d \tau_1 \int_{0}^{\tau_1}\d \tau_2 \cdots \int_{0}^{\tau_{\ell-1}} \d \tau_{\ell} \Phi(T, \tau_1) \L \Phi(\tau_1, \tau_2) \L \cdots \Phi(\tau_{\ell-1}, \tau_{\ell}) \L\Phi(\tau_{\ell},0)\B)\rho_0\Big],
\end{align*}
and \begin{align*}
    F_{T,K}(O,\lambda):=\tr\B[O\B(\int_{0}^T \d \tau_1 \int_{0}^{\tau_1}\d \tau_2 \cdots \int_{0}^{\tau_{K-1}} \d \tau_{K} \Phi(T, \tau_1) \L \Phi(\tau_1, \tau_2) \L \cdots \Phi(\tau_{K-1}, \tau_{K}) \L\tilde{\Phi}_{\lambda}(\tau_{K},0)\B)\rho_0\B].
\end{align*}
The reminder $F_{T,K}(O,\lambda)$ depends on $\lambda$ since $\tilde{\Phi}_{\lambda}(\tau_K,0)$ depends on $\lambda$.

Now, suppose that the quantum circuit $U$ itself has algorithmic error like the Trotter error. Denote the quantum circuit with a tunable step size $s$ by $U_s$. The quantum circuit $U_s$ tends to the ideal unitary $U_0$ as the step size $s\to0$. Our goal is to estimate $\tr[OU_0\rho_0U_0^{\dagger}]$ under the Markovian noise described by $\lambda\L$. A direct method is to first fix the step size $s$ and extrapolate the noise strength $\lambda$ to the $\lambda\to 0$ limit, which gives an estimate of $\tr[OU_s\rho_0 U_s^{\dagger}]$, and then extrapolate the step size $s$ to the $s\to0$ limit $\tr[OU_0\rho_0U_0^{\dagger}]$.

 What if we set $s= s(\lambda)$ depending on $\lambda$ and simultaneously extrapolate to the $s, \lambda \to 0$ limit? We first extend the previous analysis to incorporate different step sizes $s$. For different step sizes $s$, the running time $T_s$ of the algorithm can be different, and the ideal Hamiltonian $H_s(t)$, $t\in [0,T_s]$ are also different. \sz{I'm wondering if $H_s(t)$ is the effective Hamiltonian equivalent to algorithm as Trotter-Suzuki, or with additional scaling? If so, we only need the following parameters $T=sT_{\mathrm{sim}}$, $T$ and $T_s$ are actually the same thing.} \xz{No, $H_s(t)$ is the control Hamiltonian of the physical device.} We define the time-evolution operator of the ideal algorithm with step size $s$ and the noisy algorithm with noise strength $\lambda$ as \begin{align*}
    \Phi_s(t_2,t_1):=\T \exp(\int_{t_1}^{t_2} -\i\ad_{H_s(t)}\,\d t),\quad \tilde{\Phi}_{s,\lambda}(t_2, t_1):=\T \exp(\int_{t_1}^{t_2}(-\i \ad_{H_s(t)}+\lambda \L)\,\d t). 
\end{align*}
Analogous to Eq.~\eq{expansion-lambda}, we have \begin{align}
\label{eq:expansion-s-l}
    \tr[O\tilde{\Phi}_{s,\lambda}(T_s, 0)\rho_0] &=\tr[O\Phi_s(T_s, 0)\rho_0]+\sum_{\ell=1}^{K-1}F_{s,\ell}(O)\lambda^{\ell}+F_{s,K}(O,\lambda)\lambda^K \nonumber \\
    &=\tr[OU_s \rho_0 U_s^{\dagger}] +\sum_{\ell=1}^{K-1}F_{s,\ell}(O)\lambda^{\ell}+F_{s,K}(O,\lambda)\lambda^K,
\end{align}
where \begin{align}
    \label{eq:def-Fsl}
    F_{s,\ell}(O):=\tr\B[O\B(\int_{0}^{T_s} \d \tau_1 \int_{0}^{\tau_1}\d \tau_2 \cdots \int_{0}^{\tau_{\ell-1}} \d \tau_{\ell} \Phi_s(T_s, \tau_1) \L \Phi_s(\tau_1, \tau_2) \L \cdots \Phi_s(\tau_{\ell-1}, \tau_{\ell}) \L\Phi_s(\tau_{\ell},0)\B)\rho_0\Big],
\end{align}
and \begin{align*}
    F_{s, K}(O,\lambda):=\tr\B[O\B(\int_{0}^{T_s} \d \tau_1 \int_{0}^{\tau_1}\d \tau_2 \cdots \int_{0}^{\tau_{K-1}} \d \tau_{K} \Phi_s(T_s, \tau_1) \L \Phi_s(\tau_1, \tau_2) \L \cdots \Phi_s(\tau_{K-1}, \tau_{K}) \L\tilde{\Phi}_{s,\lambda}(\tau_{K},0)\B)\rho_0\B].
\end{align*}
\subsection{Take $U$ to be the whole quantum algorithm \xz{hard to proceed.}} 
To analyze the extrapolation error, we need to express Eq.~\eq{expansion-s-l} as a polynomial of $s,\lambda$ plus some reminder. For Trotter-based Hamiltonian simulation algorithms, the first term in Eq.~\eq{expansion-s-l} can be written as \[
\tr[OU_{s} \rho_{0} U_s^{\dagger}] = \tr[OU_0 \rho_0 U_0^{\dagger}]+\sum_{j=1}^{K-1} E_{T_s, j}(O)s^j+R_{K}(O,s),
\] 
\tyl{to align with the notation in Eq. (1), maybe we shall write $E_{j, K}$ as $E_{T, j}$?}\xz{Yes, and now the total physical running time $T_s$ depends on step size $s$.}
which is a polynomial of $s$ plus a reminder (see \cite{watson2024exponentially}). However, it is hard to write $F_{s,\ell}(O)$ and $F_{s,\ell}(O)$ \tyl{should be an $E$ term here?} as simple functions of $s$. The difficulty lies in expressing $\Phi_s(t_2, t_1)$ and  $\tilde{\Phi}_{s,\lambda}(t_2,t_1)$ as function of $s$ for general $t_1,t_2$. Recall that $\Phi_s(t_2, t_1)$ is the evolution of the system executing the Trotter-based Hamiltonian simulation algorithm from time $t_1$ to $t_2$. The definition of $F_{s,\ell}(O)$ (Eq.~\eq{def-Fsl}) involves $\Phi_s(t_2, t_1)$ for all $t_1,t_2\in [0,T_s]$ including $t_1, t_2$ in the middle of one Trotter step, where $\Phi_s(t_2, t_1)$ is hard to analyze. 

\subsection{Take $U$ to be one Trotter step}
Consider the Hamiltonian $H = \sum_{\gamma=1}^{\Gamma} H_{\gamma}$. Let $\P_2(t)=e^{-\i H_1 t/2}\cdots e^{-\i H_{\Gamma} t/2} e^{-\i H_{\Gamma} t/2} \cdots e^{-\i H_1 t/2}$ be the second order Trotter formula. \tyl{By writing as this, are you suggesting that $H_{1},\ldots,H_{\Gamma}$ are time-independent here?} \xz{Yes, we consider simulation of time-independent Hamiltonian first. The time-dependent Hamiltonian $H_s(t)$ in the previous parts is the control Hamiltonian of the quantum computer realizing the Hamiltonian simulation algorithm. We need the information of control Hamiltonian $H_s(t)$ of the real quantum device to analyze the impact of the physical noise $\L$. } We aim to simulate $e^{-\i H T}$ by Trotter methods. The simulation is performed by repeatedly applying $\P_2(sT)$ for $1/s$ times. \textbf{Here we assume that to implement the Trotter step for different step size, we only need to linearly scale the control Hamiltonian of the physical device.} Specifically, each application of $\P_2(sT)$ is realized by evolving the quantum computer under a time-dependent control Hamiltonian $\Hc(t/s)$ for $t\in [0, sT_{\simu}]$, where $T_{\simu}$ is the physical time required to execute one Trotter step of simulated time $T$. We take $H(t)$ to be $\Hc(t/s)$ and $T = s T_{\simu}$. Then the idea time-evolution operator is 
\begin{align*}
    \Phi_s(t_2,t_1)=\T \exp(\int_{t_1}^{t_2} -\i\ad_{\Hc(t/s)}\,\d t) = \T \exp(s\int_{t_1/s}^{t_2/s} -\i\ad_{\Hc(t)}\,\d t) 
\end{align*} 
Substituting $T_s = sT_{\simu}$ and $\tau_j=s\sigma_j$ in \eq{def-Fsl}, we get \begin{align*}
    F_{s,\ell}(O)=s^\ell \tr\left[O \int_{0}^{T_{\simu}} \d\sigma_1 \cdots \int_{0}^{\sigma_{\ell-1}} \d\sigma_{\ell} \left( \prod_{j=1}^{\ell} \Phi_s(s\sigma_{j-1}, s\sigma_j) \L \right) \Phi_s(s\sigma_\ell, 0) \rho_0 \right],
\end{align*}
where \begin{align*}
    \Phi_s(s\sigma_{j-1}, s\sigma_j) = \T \exp(s\int_{\sigma_{j}}^{\sigma_{j-1}} -\i\ad_{\Hc(t)}\,\d t)
\end{align*}
has Dyson series \begin{align*}
    \Phi_s(s\sigma_{j-1}, s\sigma_j)=\sum_{k=0}^{\infty} s^k V_k,
\end{align*}
for \begin{align*}
    \left\|V_k\right\| \leq \frac{\left(2 H_{\max }\left(\sigma_{j-1}-\sigma_j\right)\right)^k}{k!}.
\end{align*}
\bibliographystyle{plainnat}
\bibliography{ref}
\end{document}
